\subsection{Sprint 4}

Chegando à Sprint 4, iniciamos a planning avaliando o que havia sido feito até então,
o que precisava ser finalizado das sprints anteriores e o que ainda faltava entregar.
A conclusão foi \textbf{excluir as funcionalidades do usuário professor}, pois não
haveria tempo hábil para implementá-las. Essa decisão já havia sido comunicada à
stakeholder, que compreendeu e aceitou tranquilamente, ciente de que a construção
desse perfil poderia ser continuada futuramente.

Nesta sprint, realizei, com apoio de colegas, \textbf{três tasks}:

\subsubsection*{Front-end --- Issue 82: Cadastrar usuário aluno}

No front-end, trabalhei na Issue 82, referente ao cadastro de usuário aluno. Tive
dificuldade em definir a melhor abordagem para a task, pois adicionar um novo ícone
à interface poluiria a área com informações em excesso. A ideia foi criar, dentro do
\textit{``editar aluno''}, uma opção para adicionar senha e criar o usuário; após o
usuário ser criado, essa opção não apareceria mais.

Comecei a implementar a integração, mas não conseguia conectar com o backend por
falta de prática com front-end. Pedi ajuda ao Mauro, que se disponibilizou para uma
call na sexta-feira, 21/06, às 21h. Começamos a implementação pelo backend, mas
encontramos diversos erros --- alguns \textit{endpoints} não retornavam ao front-end
o que era solicitado. Às 23h30, fizemos o push no estado em que estava para tentar
resolver no dia seguinte. Encaminhei o código para a Laura, que já havia tentado a
mesma abordagem. Ela retornou com o \textit{feedback} de que também encontrou
dificuldades com nossa abordagem e preferiu criar um \textit{modal} simples com um
botão para criar o usuário, conforme a proposta inicial, pois era mais prático. Ela
entregou o \ac{mr} funcionando e testado.

\subsubsection*{Back-end --- Issue 23: Adicionar professor à aula}

A Issue 23 consistia em adicionar o professor à aula. Identificamos que a entidade
\textit{course} não possuía um professor cadastrado como responsável principal.
As mudanças foram objetivas: adicionamos o campo professor na \textit{Entity course}
e ajustamos as demais entidades relacionadas, incluindo \textit{controller},
\textit{service} e DTOs com as novas \textit{features}.

\subsubsection*{Back-end --- Issue 24: Adicionar aluno temporário à aula}

A Issue 24 era adicionar aluno temporário à aula. Inicialmente imaginei que seria
similar à Issue 23, mas durante o desenvolvimento encontrei muitos conflitos entre
as entidades. Preferi pedir ajuda ao Mauro, que se disponibilizou novamente para uma
call na sexta-feira à noite. Sua contribuição foi essencial: ele propôs trabalhar de
forma a não alterar muito a entidade \textit{course}, recebendo uma lista de alunos e
centralizando a lógica na entidade \textit{Aluno}. Implementamos uma série de
pequenas alterações e validações no sistema, testamos os inserts e tudo funcionou
corretamente. Aproveitei também para pedir que ele revisasse a Issue 23, e ele deu
dicas valiosas para otimizar a implementação. Foi um dia extremamente produtivo.
